% =====================================================================
%  Chapter 21 — Gravity
% =====================================================================

\chapter{Gravity: The Detection We Already Have}
\label{ch:gravity}

The oldest sense is also the one the field is least inclined to count, because
it is old. The argument of this chapter is that it should be counted, because
what the field has achieved with gravity is unmatched by anything it has
attempted with a detector.

\section{What has been achieved}In ninety years of dark matter physics, the
gravitational use of dark matter has produced a set of measurements that no
other technique in cosmology comes close to.

The cosmic microwave background is the centrepiece. A 2.7~kelvin
temperature field, uniform to one part in a hundred thousand, has been mapped
by four independent instruments to a precision that resolves the acoustic
peaks of the photon-baryon fluid. The positions of those peaks measure the
dark matter density to two per cent; their relative heights measure the baryon
density to one per cent; their damping tail constrains the free-streaming
scale; and the lensing in the same maps measures the growth of structure
between recombination and the present. That is four measurements, from one
field, of four different aspects of dark matter.

Gravitational lensing has turned the same object into a probe of the
distribution of matter at scales from the kiloparsec to the megaparsec.
Weak lensing maps the cosmic web and its mass; strong lensing measures
individual haloes' mass with no luminosity assumption; and the combination of
lensing with spectroscopic redshifts from DESI gives a three-dimensional map
that has become one of the main competitors to the CMB for structure
constraints.

And the individual systems remain indispensable. The Bullet Cluster measures
the collisionless property of the dark component. The rotation curves measure
the radial distribution. The dwarf galaxies measure the low-mass end. Each
of these is a measurement that would be impossible without dark matter, and
each is a measurement of something different.

\section{What gravity cannot measure}The list is short and it is the list of
what the field actually wants to know.

Gravity measures \emph{abundance}: how much, where, and at what scale. It
does not measure \emph{mass}, because a particle of any mass produces the
same gravitational effect. It does not measure \emph{couplings}, because the
gravitational coupling is universal. It does not measure \emph{spin}, or
whether the dark matter is a particle, a field, or a population of black
holes. And it cannot detect an individual dark matter particle, because the
gravitational effect of one particle at the Earth's distance is
unimaginably below any measurable quantity.

This is the precise sense in which the gravitational programme is finished.
It has measured everything it can measure, and the things it cannot measure
require a different kind of experiment entirely.

\section{The conversion argument}There is an argument worth making explicitly,
because it is the strongest answer to the question of what dark matter
research has achieved.

The gravitational measurements are not a preliminary to the real work. They
are the specification. The dark matter density they establish is the number
that every candidate must reproduce; the free-streaming scale they constrain
is the number that every production mechanism must evade; the small-scale
structure they leave unexplained is the number that every model must
explain. The laboratory programme is a programme of testing a
specification, and the specification is a product of the gravitational
programme.

Any assessment of dark matter research that counts the gravitational results
as preliminary and the laboratory results as primary has the history
backwards.

\begin{keytakeaways}
\item The gravitational programme has produced four distinct measurements
from the cosmic microwave background alone, plus lensing maps and a set of
object-specific measurements that no other technique can make.
\item It cannot measure mass, coupling, spin or identity, and it cannot
detect an individual particle, so it is finished as a measurement programme.
\item The gravitational results are the specification that every candidate
must satisfy, not a preliminary to the real work.
\end{keytakeaways}
